% ============================================================================
%  Systematic MVAR vs LSTM Comparison for Crowd Density ROM
% ============================================================================
\documentclass[11pt]{article}

% ---- Packages ----
\usepackage[margin=1in]{geometry}
\usepackage{amsmath, amssymb, amsthm}
\usepackage{booktabs}
\usepackage{multirow}
\usepackage{hyperref}
\usepackage{xcolor}
\usepackage{enumitem}
\usepackage[font=small]{caption}
\usepackage{subcaption}
\usepackage{float}
\usepackage{array}
\usepackage{colortbl}
\usepackage{longtable}
\usepackage{tikz}
\usepackage{pgfplots}
\pgfplotsset{compat=1.18}
\usepackage{siunitx}

% ---- Colors ----
\definecolor{mvarblue}{HTML}{1565C0}
\definecolor{lstmorange}{HTML}{E65100}
\definecolor{goodgreen}{HTML}{2E7D32}
\definecolor{badred}{HTML}{C62828}
\definecolor{accentblue}{HTML}{1565C0}
\definecolor{lightgray}{HTML}{F5F5F5}

% ---- Custom commands ----
\newcommand{\R}{\mathbb{R}}
\newcommand{\bx}{\mathbf{x}}
\newcommand{\ba}{\mathbf{a}}
\newcommand{\bU}{\mathbf{U}}
\newcommand{\bA}{\mathbf{A}}
\newcommand{\MVAR}{\textcolor{mvarblue}{\textbf{MVAR}}}
\newcommand{\LSTM}{\textcolor{lstmorange}{\textbf{LSTM}}}

\hypersetup{
    colorlinks=true,
    linkcolor=accentblue,
    citecolor=accentblue,
    urlcolor=accentblue,
}

\begin{document}

% ============================================================================
\section{Mathematical Formulation}
\label{sec:math}
% ============================================================================

\subsection{Particle Dynamics}

We simulate $N=300$ self-propelled particles in a periodic
$[0, L_x]\times[0, L_y]$ domain ($L_x = L_y = 25$).  Each particle $i$
has position $\bx_i\in\R^2$ and heading $\theta_i\in[0,2\pi)$.

\paragraph{Vicsek alignment.}
At each timestep, particle $i$ averages the headings of neighbors within
radius $R=4$:
\begin{equation}
    \theta_i^{(t+1)} = \arg\Bigl(\sum_{j:\|\bx_j-\bx_i\|<R}
        e^{i\theta_j^{(t)}}\Bigr) + \eta\,\xi_i^{(t)},
    \qquad \xi_i \sim \mathrm{Uniform}(-\pi,\pi),
    \label{eq:vicsek}
\end{equation}
where $\eta\in[0.03, 1.5]$ is the noise intensity.

\paragraph{Morse forces.}
An additional pairwise Morse-type force modifies particle velocities:
\begin{equation}
    \mathbf{F}_{ij} = -\nabla U_{\text{Morse}}(r_{ij}),
    \qquad
    U_{\text{Morse}}(r) = C_r\, e^{-r/\ell_r} - C_a\, e^{-r/\ell_a},
    \label{eq:morse}
\end{equation}
where $C_a, C_r$ are attraction/repulsion strengths and $\ell_a, \ell_r$ the
corresponding length scales.  The net force is truncated at
$r_{\text{cut}} = 5\,\ell_a$.

\paragraph{Position update.}
\begin{equation}
    \bx_i^{(t+1)} = \bx_i^{(t)} + v_0\,
    \begin{pmatrix}\cos\theta_i^{(t+1)}\\\sin\theta_i^{(t+1)}\end{pmatrix}
    \Delta t + \mu_t\,\mathbf{F}_i^{(t)}\,\Delta t,
    \label{eq:position}
\end{equation}
with self-propulsion speed $v_0$, force coupling $\mu_t=0.3$, and
$\Delta t=0.04$\,s.  In ``variable speed'' variants, $v_0$ is modulated
by the local force magnitude.

\subsection{Kernel Density Estimation}

Particle positions are converted to a continuous density field via a
Gaussian KDE on a $64\times 64$ grid:
\begin{equation}
    \rho(\mathbf{r}) = \frac{1}{N}\sum_{i=1}^N
        \frac{1}{2\pi h^2}\exp\!\Bigl(-\frac{\|\mathbf{r}-\bx_i\|^2}{2h^2}\Bigr),
    \label{eq:kde}
\end{equation}
with bandwidth $h=5.0$.  Each snapshot is flattened to a vector
$\rho\in\R^{4096}$.

\subsection{Proper Orthogonal Decomposition (POD)}

Let $\mathbf{S}\in\R^{4096\times (M\cdot T_\mathrm{rom})}$ be the
snapshot matrix (all training runs, all subsampled timesteps).  We
perform shift alignment (subtract per-run spatial mean) and then compute
the truncated SVD:
\begin{equation}
    \mathbf{S} - \bar{\mathbf{S}}
    \approx \bU\,\boldsymbol{\Sigma}\,\mathbf{V}^\top,
    \qquad
    \bU\in\R^{4096\times R},\quad R=19.
    \label{eq:pod}
\end{equation}
The number of modes $R=19$ is fixed across all experiments.  The latent
representation of any density snapshot is:
\begin{equation}
    \ba(t) = \bU^\top\bigl(\rho(t) - \bar\rho\bigr) \in \R^{19}.
    \label{eq:project}
\end{equation}

\subsection{Model 1: MVAR (Multivariate AutoRegressive)}
\label{sec:mvar}

The MVAR model assumes a linear relationship between past and future
latent states.  Given lag $p=5$:
\begin{equation}
    \ba(t) = \sum_{k=1}^{p} \bA_k\, \ba(t-k) + \mathbf{b} + \boldsymbol{\epsilon},
    \label{eq:mvar}
\end{equation}
where $\bA_k\in\R^{19\times 19}$ are coefficient matrices and
$\mathbf{b}\in\R^{19}$ is a bias.  All $\{\bA_k, \mathbf{b}\}$ are
trained jointly via Ridge regression:
\begin{equation}
    \min_{\bA,\mathbf{b}}\;\sum_{t=p+1}^{T}
    \Bigl\|\ba(t) - \sum_{k=1}^p \bA_k\,\ba(t-k) - \mathbf{b}\Bigr\|^2
    + \alpha\sum_{k=1}^p\|\bA_k\|_F^2,
    \qquad \alpha=10^{-4}.
    \label{eq:ridge}
\end{equation}

\paragraph{Parameter count.}
$p\times R\times R + R = 5\times19\times19 + 19 = \mathbf{1{,}824}$.

\paragraph{Stability analysis.}
The companion matrix $\mathbf{C}\in\R^{pR\times pR}$ is:
\begin{equation}
    \mathbf{C} = \begin{pmatrix}
        \bA_1 & \bA_2 & \cdots & \bA_{p-1} & \bA_p \\
        \mathbf{I} & \mathbf{0} & \cdots & \mathbf{0} & \mathbf{0}\\
        \mathbf{0} & \mathbf{I} & \cdots & \mathbf{0} & \mathbf{0}\\
        \vdots & & \ddots & & \vdots\\
        \mathbf{0} & \mathbf{0} & \cdots & \mathbf{I} & \mathbf{0}
    \end{pmatrix}.
    \label{eq:companion}
\end{equation}
The spectral radius $\rho(\mathbf{C})=\max_i|\lambda_i|$ governs
long-horizon stability: if $\rho>1$, forecasts eventually diverge.
Our pipeline logs $\rho$ and optionally enforces stability via uniform
coefficient scaling.

\paragraph{Forecasting.}
At test time, the model autoregressively rolls forward:
\begin{equation}
    \hat{\ba}(t) = \sum_{k=1}^p \bA_k\,\hat\ba(t-k) + \mathbf{b},
    \qquad t = p+1,\ldots,T_{\text{test}}/\Delta t.
    \label{eq:mvar_forecast}
\end{equation}

\subsection{Model 2: LSTM (Long Short-Term Memory)}
\label{sec:lstm}

The LSTM learns a nonlinear map from a window of past latent states
to the next state:
\begin{equation}
    \hat\ba(t) = f_\theta\!\bigl(\ba(t-L),\ldots,\ba(t-1)\bigr),
    \qquad L = 20,
    \label{eq:lstm_map}
\end{equation}
where $f_\theta$ is a 2-layer stacked LSTM with 128 hidden units per
layer.

\paragraph{Architecture details.}
\begin{itemize}[nosep]
    \item Input: $[\text{batch}, L, 19]$ sequence of latent vectors
        (z-scored per mode using training-set statistics).
    \item LSTM: 2 layers, $N_h=128$ hidden units, dropout 0.1 between layers.
    \item LayerNorm on last hidden state $\mathbf{h}_{\text{last}}\in\R^{128}$.
    \item Output head: $\mathbf{W}\mathbf{h}_{\text{last}}+\mathbf{b}\in\R^{19}$.
    \item Residual connection:
        $\hat\ba(t) = \ba(t-1) + \Delta\hat\ba(t)$
        (network predicts the increment $\Delta$).
\end{itemize}

\paragraph{Parameter count.}
\begin{align}
    \text{Layer 1:}\quad & 4(N_h \cdot d + N_h^2 + 2N_h) = 4(128\cdot19+128^2+256) = 76{,}800,\notag\\
    \text{Layer 2:}\quad & 4(N_h^2 + N_h^2 + 2N_h) = 4(128^2+128^2+256) = 132{,}096\cdot\tfrac{1}{2}\approx 66{,}048,\notag\\[-2pt]
    \text{LayerNorm:}\quad & 2N_h = 256,\notag\\
    \text{Output:}\quad & N_h\cdot d + d = 128\cdot19+19 = 2{,}451,\\
    \text{Total:}\quad & \approx\mathbf{106{,}000}.\notag
    \label{eq:lstm_params}
\end{align}

\paragraph{Training protocol.}
\begin{itemize}[nosep]
    \item Adam optimizer with learning rate $10^{-3}$, weight decay $10^{-5}$.
    \item Gradient clipping at norm 1.0.
    \item 80/20 train/validation split; early stopping with patience 30.
    \item Maximum 300 epochs; batch size 256.
    \item MSE loss in normalized latent coordinates.
\end{itemize}

\subsection{Lifting and Evaluation}

The predicted latent trajectory $\hat\ba(t)$ is lifted back to the
physical density field:
\begin{equation}
    \hat\rho(t) = \bU\,\hat\ba(t) + \bar\rho.
    \label{eq:lift}
\end{equation}

Quality is measured by the coefficient of determination and RMSE:
\begin{equation}
    R^2(t) = 1 - \frac{\|\rho(t) - \hat\rho(t)\|^2}
                       {\|\rho(t) - \bar\rho_{\text{test}}\|^2},
    \qquad
    \text{RMSE}(t) = \|\rho(t) - \hat\rho(t)\|_2.
    \label{eq:metrics}
\end{equation}
We report both a per-timestep $R^2(t)$ curve and the time-averaged
overall $R^2$.  The ``POD ceiling'' $R^2_{\text{POD}}$ (projection error
of the true trajectory onto the basis) provides an upper bound on
forecast quality.


% ============================================================================
\section{Experimental Design: 54 Physical Configurations}
\label{sec:configs}
% ============================================================================

All 54 configs are generated by \texttt{generate\_regime\_configs.py} and
share identical simulation parameters ($N=300$, $L=25$, $\Delta t=0.04$,
$T_{\text{train}}=20$\,s, $T_{\text{test}}=200$\,s), KDE settings
($64\times64$, $h=5.0$), POD settings ($R=19$, shift alignment), and ROM
hyperparameters (MVAR: $p=5$, $\alpha=10^{-4}$; LSTM: $L=20$,
$N_h=128$, 2 layers).  Only the \emph{physics} varies.

\vspace{0.5em}
The 54 experiments decompose as:
\begin{itemize}[nosep]
    \item \textbf{NDYN suite} (14~regimes): Vicsek--Morse with extreme diversity.
    \item \textbf{d'Orsogna suite} (14~regimes): $C_r/C_a$ vs $\ell_r/\ell_a$ taxonomy.
    \item Each regime $\times$ 2 speed modes (constant, variable) $-$ exceptions.
\end{itemize}

\subsection{NDYN Suite: 14 Diverse Vicsek--Morse Regimes}

Each NDYN configuration occupies a distinct corner of the physical
parameter space.  At least two parameter axes differ between any pair.

\begin{longtable}{llp{1.5cm}p{1.0cm}p{5.5cm}}
    \toprule
    \textbf{ID} & \textbf{Name} & \textbf{Speed, $\eta$} & \textbf{$C_a$/$C_r$} & \textbf{Description} \\
    \midrule
    \endfirsthead
    \toprule
    \textbf{ID} & \textbf{Name} & \textbf{Speed, $\eta$} & \textbf{$C_a$/$C_r$} & \textbf{Description} \\
    \midrule
    \endhead

    01 & Crawl        & 0.5, 0.03 & 0.8/0.3 & Slow deterministic flock; easy POD baseline \\
    02 & Flock        & 3.0, 0.10 & 2.0/0.5 & Ordered coherent traveling cluster \\
    03 & Sprint       & 10.,\;0.10 & 2.0/0.5 & Fast transport; tests shift alignment \\
    04 & Gas          & 3.0, 1.50 & 0.5/0.2 & High noise destroys all structure \\
    05 & Black Hole   & 4.0, 0.15 & 20./0.05 & Extreme attraction $\to$ density spike \\
    06 & Supernova    & 3.0, 0.15 & 0.05/15. & Explosive repulsion; violent dispersal \\
    07 & Crystal      & 2.0, 0.10 & 3.0/3.0 & Balanced forces $\to$ ring/lattice equilibrium \\
    08 & Pure Vicsek  & 3.0, 0.30 & \emph{none} & No forces; alignment-only band patterns \\
    09 & Long Range   & 2.0, 0.20 & 1.5/0.6 & $\ell_a{=}6$, $\ell_r{=}3$; global coupling \\
    10 & Short Range  & 3.0, 0.15 & 4.0/2.0 & $\ell_a{=}0.3$, $\ell_r{=}0.1$; clumpy fragments \\
    11 & Noisy Collapse & 3.0, 0.80 & 12./0.1 & Strong attraction vs high noise \\
    12 & Fast Explosion & 8.0, 0.20 & 0.1/10. & High speed + strong repulsion \\
    13 & Chaos        & 12.,\;1.20 & 1.0/0.5 & Extreme speed+noise; ROM lower bound \\
    14 & Var.\,Speed  & 3.0, 0.20 & 2.0/0.8 & Speed coupled to forces \\
    \bottomrule
\end{longtable}

All except NDYN08 (no forces) and NDYN14 (already variable-speed)
get a \texttt{\_VS} duplicate with \texttt{speed\_mode=variable},
yielding $14 + 12 = 26$ NDYN experiments.

\subsection{d'Orsogna Suite: 14 Regimes from the $(C,\ell)$ Taxonomy}

Following Bhaskar \& Ziegelmeier, we parametrize d'Orsogna dynamics by
the ratios $C \equiv C_r/C_a$ and $\ell \equiv \ell_r/\ell_a$, fixing
$C_a=1$, $\ell_a=1$ as reference scales.

\begin{longtable}{llccp{5.5cm}}
    \toprule
    \textbf{ID} & \textbf{Regime Class} & $C$ & $\ell$ & \textbf{Expected Behavior} \\
    \midrule
    \endfirsthead
    \toprule
    \textbf{ID} & \textbf{Regime Class} & $C$ & $\ell$ & \textbf{Expected Behavior} \\
    \midrule
    \endhead

    SM01 & Single Mill        & 0.5 & 0.1 & Rotational vortex (weak $C$) \\
    SM02 & Single Mill        & 3.0 & 0.1 & Rotational vortex (strong $C$) \\
    SM03 & Single Mill        & 2.0 & 0.5 & Strong $C$, medium range \\
    DM01 & Double Mill        & 0.9 & 0.5 & Counter-rotating double vortex \\
    DR01 & Double Ring        & 0.1 & 0.1 & Nested annular structures (small) \\
    DR02 & Double Ring        & 0.9 & 0.9 & Nested annular structures (large) \\
    CS01 & Collective Swarm   & 0.1 & 0.5 & Cohesive traveling group (low $C$) \\
    CS02 & Collective Swarm   & 0.5 & 3.0 & Cohesive traveling group (med $C$) \\
    CS03 & Collective Swarm   & 0.9 & 3.0 & Cohesive traveling group (high $C$) \\
    ES01 & Escape (sym.)      & 3.0 & 0.9 & Symmetric particle escape \\
    EU01 & Escape (unsym.)    & 2.0 & 2.0 & Asymmetric dispersal \\
    EU02 & Escape (unsym.)    & 3.0 & 3.0 & Asymmetric dispersal (extreme) \\
    EC01 & Escape (collective) & 2.0 & 3.0 & Group escape \\
    EC02 & Escape (collective) & 3.0 & 0.5 & Group escape (short range) \\
    \bottomrule
\end{longtable}

All 14 receive a \texttt{\_VS} variable-speed variant, yielding $14 + 14
= 28$ d'Orsogna experiments.  Combined with the 26 NDYN experiments, the
total is $\mathbf{54}$.


% ============================================================================
\section{Initial Conditions}
\label{sec:ics}
% ============================================================================

Diversity of initial conditions is critical to avoid overfitting to
a single spatial configuration.  All 54 experiments use the same IC
distributions:

\subsection{Training ICs (168 total per experiment)}

\begin{longtable}{lcp{8cm}}
    \toprule
    \textbf{IC Type} & \textbf{Count} & \textbf{Generation Strategy} \\
    \midrule
    Gaussian      & 80 & Centers on $4\times4$ grid $\{5,10,15,20\}^2$;
                         variances $\sigma^2\in\{1.5, 3.0, 5.0\}$;
                         1 sample per configuration.\\[3pt]
    Uniform       & 40 & Particles drawn uniformly from $[0,L]^2$. \\[3pt]
    Two-Cluster   & 24 & Two Gaussian clusters with separations
                         $d\in\{4,7,10\}$, widths $\sigma\in\{1.5,3.0\}$;
                         4 samples per config. \\[3pt]
    Ring          & 24 & Particles on annulus with radii
                         $r\in\{3,5,8\}$, widths $w\in\{0.5,1.0\}$;
                         4 samples per config. \\
    \bottomrule
\end{longtable}

\subsection{Test ICs (4 total per experiment)}

\begin{longtable}{lcp{9cm}}
    \toprule
    \textbf{IC Type} & \textbf{Count} & \textbf{Parameters} \\
    \midrule
    Gaussian      & 1 & Center $(12.5, 12.5)$, $\sigma^2=2.5$ \\
    Uniform       & 1 & Full-domain uniform \\
    Two-Cluster   & 1 & Separation 7, width 2.0 \\
    Ring          & 1 & Radius 5.0, width 0.8 \\
    \bottomrule
\end{longtable}

Test parameters are chosen to lie between (not on) the training grid
values, ensuring genuine interpolation tests.


% ============================================================================
\section{Shared POD Basis and Latent Dataset}
\label{sec:pod}
% ============================================================================

For each experiment, the 168 training simulations produce density
snapshots that are subsampled by a factor of 3 (every 3rd timestep at
$\Delta t_{\text{rom}} = 0.12$\,s).  After shift alignment (subtracting
each run's spatial mean), all snapshots are stacked into the global
snapshot matrix $\mathbf{S}\in\R^{4096\times (168 \cdot T_{\text{rom}})}$ and
the top $R=19$ left singular vectors form the shared POD basis.

Both models operate on the same latent dataset
$\{\ba_m(t)\}_{m=1}^{168}$, where
$\ba_m(t) = \bU^\top(\rho_m(t) - \bar\rho)$.

Key outputs at this stage:
\begin{itemize}[nosep]
    \item \texttt{rom\_common/pod\_basis.npz}: $\bU$, $\bar\rho$, singular values, energy spectrum.
    \item \texttt{rom\_common/latent\_trajectories.npz}: Complete latent dataset.
\end{itemize}


% ============================================================================
\section{Visualization and Analysis Pipeline}
\label{sec:viz}
% ============================================================================

After Oscar produces test artifacts, a local 10-step visualization
pipeline (\texttt{run\_visualizations.py}) generates all plots and
reports.  The pipeline auto-discovers which models (MVAR, LSTM) have
results on disk, so it adapts gracefully if one model is absent.

\subsection{Step 1: POD Diagnostic Plots}

\begin{itemize}[nosep]
    \item \texttt{pod\_singular\_values.png}: Normalized singular value spectrum $\sigma_i/\sigma_1$.
    \item \texttt{pod\_energy.png}: Cumulative energy $\sum_{i=1}^k \sigma_i^2 / \sum_i \sigma_i^2$.
    \item \texttt{pod\_modes\_grid.png}: First 6 spatial modes $\bU_{:,i}$ reshaped to $64\times64$.
\end{itemize}

\subsection{Step 2: Test Metrics Computation}

For each model and each of the 4 test ICs:
\begin{itemize}[nosep]
    \item Load true density $\rho(t)$ and predicted trajectory $\hat\rho(t)$.
    \item Compute time-resolved $R^2(t)$, $\text{RMSE}(t)$, and overall scalars.
    \item Save \texttt{test\_results.csv} with columns: \texttt{test\_id},
        \texttt{ic\_type}, \texttt{r2\_reconstructed}, \texttt{r2\_latent},
        \texttt{r2\_pod}, \texttt{rmse}, etc.
\end{itemize}

\subsection{Step 3: Best-Run Visualizations}

For the test run with the highest $R^2$:
\begin{itemize}[nosep]
    \item Side-by-side density snapshots (true vs predicted) at $t=0, T/4, T/2, 3T/4, T$.
    \item Latent trajectory comparison in the top-3 POD modes.
    \item Pointwise error heatmaps $|\rho - \hat\rho|$.
\end{itemize}

\subsection{Step 4: Summary Plots}

These are the core comparison plots:
\begin{itemize}[nosep]
    \item \texttt{r2\_by\_ic\_type.png}: Grouped bar chart of $\overline{R^2}$ by IC type,
        one bar per model.
    \item \texttt{error\_by\_ic\_type.png}: Grouped bar chart of $\overline{\text{RMSE}}$ by IC type.
    \item Per-test scatter showing $R^2_{\text{lifted}}$ vs $R^2_{\text{latent}}$
        for both models.
\end{itemize}

\subsection{Step 5: Time-Resolved Analysis}

Captures long-horizon degradation behavior:
\begin{itemize}[nosep]
    \item \texttt{r2\_mean\_over\_time.png}: $\overline{R^2}(t)$ curve for each model,
        with IQR shading. Reveals when/if $R^2$ drops below 0.
    \item \texttt{survival\_curves.png}: Fraction of tests with $R^2(t)>0.5$
        and $R^2(t)>0$ vs time.
    \item \texttt{rmse\_over\_time.png}: Error growth over time.
    \item \texttt{r2\_heatmap\_by\_test.png}: $R^2(t)$ heatmap (tests $\times$ time).
\end{itemize}

\subsection{Step 6: LSTM Training Diagnostics}

If the LSTM was trained:
\begin{itemize}[nosep]
    \item Training/validation loss curves (log scale) with best-epoch marker.
    \item Loss improvement rate $-\Delta L_{\text{val}}$ per epoch.
\end{itemize}

\subsection{Step 7: MVAR vs LSTM 4-Panel Comparison}

A single composite figure with:
\begin{enumerate}[nosep]
    \item $R^2$ by IC type (grouped bars).
    \item RMSE by IC type (grouped bars).
    \item $R^2$ distribution (violin plot, all tests).
    \item Numerical summary table (mean, std, winner).
\end{enumerate}

\subsection{Step 8: Latent vs Lifted Analysis}

Per model:
\begin{itemize}[nosep]
    \item Time-resolved $R^2_{\text{latent}}(t)$ vs $R^2_{\text{lifted}}(t)$
        with IQR envelopes and POD ceiling.
    \item Lifting gap bar: quantifies how much $R^2$ is lost in the
        $\bU\hat\ba + \bar\rho$ reconstruction.
    \item Spectral radius annotation (MVAR only).
\end{itemize}

\subsection{Step 9: Density Videos}

For each test IC and each model: MP4 animations of true vs predicted
density fields side-by-side, with $R^2(t)$ overlay.

\subsection{Step 10: Summary JSON}

Machine-readable summary: \texttt{summary.json} with per-model mean
$R^2$, RMSE, best/worst tests, runtime, spectral radius, IC-type
breakdown, and the overall winner (MVAR or LSTM).


% ============================================================================
\section{Computational Setup}
\label{sec:compute}
% ============================================================================

\begin{itemize}[nosep]
    \item \textbf{Cluster:} Brown University Oscar HPC.
    \item \textbf{Job array:} \texttt{sbatch --array=0-53} dispatches all 54 configs.
    \item \textbf{Resources per task:} 8 cores, 48\,GB RAM, 15\,h walltime, 1 GPU.
    \item \textbf{Training sims:} 168 runs $\times$ 20\,s $\times$ 500 steps (dominant cost).
    \item \textbf{Test sims:} 4 runs $\times$ 200\,s $\times$ 5\,000 steps.
    \item \textbf{MVAR training:} $<1$\,s (Ridge regression in closed form).
    \item \textbf{LSTM training:} 2--10\,min (300 epochs, early stopping).
    \item \textbf{Estimated total:} $\sim$54 $\times$ 3\,h $\approx$ 162\,GPU-hours.
\end{itemize}


% ============================================================================
\section{Expected Outputs and Cross-Experiment Aggregation}
\label{sec:outputs}
% ============================================================================

Each of the 54 experiments produces an independent
\texttt{oscar\_output/<experiment>/} directory with:

\begin{itemize}[nosep]
    \item \texttt{rom\_common/}: POD basis, singular values, latent dataset.
    \item \texttt{MVAR/}: Model weights, test results CSV, forecasts NPZ.
    \item \texttt{LSTM/}: Model checkpoint \texttt{.pt}, training log, normalization stats, test results.
    \item \texttt{test/}: Per-run true densities, metadata JSON.
    \item \texttt{plots/}: All generated visualizations (Steps 1--10).
\end{itemize}

After downloading results, a cross-experiment summary can aggregate the
54 experiments into global rankings: which regimes are easy/hard, how
MVAR and LSTM compare by regime class, and whether variable-speed
consistently helps or hurts.


% ============================================================================
\section{Next Steps: PDE Discovery via Weak SINDy}
\label{sec:wsindy}
% ============================================================================

A natural extension of this pipeline is \emph{physics-informed}
forecasting: instead of learning a generic dynamical system in latent
space, discover the governing PDE of the density field directly.  Weak
SINDy (Sparse Identification of Nonlinear Dynamics using weak
formulations) is a promising approach for this.

\paragraph{Concept.}
Weak SINDy integrates the density field $\rho(\mathbf{r}, t)$ against
compactly-supported test functions $\psi(\mathbf{r}, t)$, converting
PDE residuals into a sparse linear system.  The sparsest model
consistent with the data is selected via an $\ell_1$ or thresholded
regression.

\paragraph{Potential benefits.}
\begin{itemize}[nosep]
    \item \emph{Interpretability:} The discovered PDE terms (advection,
        diffusion, nonlinear self-interaction) reveal the macroscopic
        physics governing the crowd.
    \item \emph{Generalization:} A physics-based model may extrapolate
        to regimes not seen during training.
    \item \emph{Complementarity:} WSINDy forecasts in the full density
        space, bypassing POD truncation error entirely.
\end{itemize}

\paragraph{Integration plan.}
Adding WSINDy as a third model would extend Steps 4--6 of the pipeline.
The same test ICs and metrics would be used, enabling a direct three-way
comparison: MVAR (linear latent) vs LSTM (nonlinear latent) vs WSINDy
(PDE in physical space).  The infrastructure for WSINDy discovery and
forecasting is being developed and will be incorporated in a future
iteration of this pipeline.


% ============================================================================
\section{Conclusion}
\label{sec:conclusion}
% ============================================================================

This document describes a systematic, controlled comparison of two
latent-space forecasting paradigms for crowd density ROM.  By fixing
every aspect of the pipeline except the forecasting model (MVAR vs LSTM)
and the underlying physics (54 regimes), we isolate the effect of model
capacity on prediction quality:

\begin{itemize}[nosep]
    \item \MVAR{}: 1{,}824 parameters, closed-form training, guaranteed
        spectral radius analysis---the minimal baseline.
    \item \LSTM{}: $\sim$106{,}000 parameters, gradient-based training
        with residual connections and layer normalization---the nonlinear
        alternative.
\end{itemize}

The 54 experiments span easy regimes (slow, deterministic flocks) to
hard regimes (fast, chaotic dispersal), with both Vicsek--Morse and
d'Orsogna force models, constant and variable speed.  The comprehensive
visualization pipeline produces 20+ plots per experiment and
machine-readable summaries, enabling both qualitative inspection and
large-scale statistical analysis.

Results from the Oscar cluster will reveal whether the LSTM's added
capacity provides consistent gains across all regimes, or whether the
simpler MVAR suffices for the quasi-linear dynamics typical of
POD-compressed crowd densities.

\end{document}
